\section{Congruence, normalizers, and discretizations}\label{part:B}

We use the conventions of Part A.  Thus $P_{g,n}$ is the pure,
orientation-preserving mapping class group, $M^+_{g,n}$ allows permutations
of the punctures, and $M^{\pm}_{g,n}$ also allows reversal of orientation.
We write $\Pi_{g,n}$ for the discrete surface group and
$\pi_{g,n}=\widehat{\Pi}_{g,n}$ for its full profinite completion.  Unless
an exception is explicitly stated, $2g-2+n>0$.  All automorphisms of
profinite groups are continuous.  In
\[
 O_{g,n}=\operatorname{Out}_{\mathrm{lac}}(\pi_{g,n})
\]
the peripheral conjugacy classes may be permuted and their generators
may be raised to a common profinite unit.  For a proper surface the
corresponding multiplier is its action on the orientation class.  We set
\[
 \check P_{g,n}=\overline{P_{g,n}},\qquad
 \Sigma_{g,n}=\overline{M^+_{g,n}},\qquad
 H_{g,n}=\overline{M^{\pm}_{g,n}},
 \qquad N^{\mathrm{norm}}_{g,n}=N_{O_{g,n}}(\Sigma_{g,n}).
\]
These are images or closures in $O_{g,n}$, not, by definition, abstract
profinite completions.  This distinction is indispensable in reading
\S26, folios 162--165, scans 320--326.  The arithmetic image
$N^{\mathrm{arith}}$ introduced in \S29, folio 188, scan 371 is a
different group; its behavior under filling is treated in Part C.

\subsection{The congruence problem and the size of the ambient group}

\statement{B.1}{Proved}{The precise congruence question}
There is a canonical continuous epimorphism
\[
 c_{g,n}:\widehat P_{g,n}\longrightarrow\check P_{g,n},
 \qquad C_{g,n}=\ker(c_{g,n}).
\]
The following assertions are equivalent: $C_{g,n}=1$; the profinite and
congruence topologies on $P_{g,n}$ coincide; and every subgroup of finite
index in $P_{g,n}$ contains the kernel of
$P_{g,n}\to\operatorname{Out}(\Pi_{g,n}/K)$ for some characteristic
subgroup $K$ of finite index in $\Pi_{g,n}$.  None of these assertions
asserts normality of $\check P_{g,n}$, $\Sigma_{g,n}$, or $H_{g,n}$ in
$O_{g,n}$.

\begin{proof}
The kernels in the statement form a neighborhood basis for the induced
congruence topology.  Indeed, a finitely generated group has only finitely
many subgroups of a prescribed finite index; intersecting all subgroups
of index at most a prescribed integer gives cofinal characteristic
subgroups.  The completion for this topology is exactly the closure of
the image in $\operatorname{Out}(\pi_{g,n})$, since the latter has the
topology determined by its characteristic finite quotients.  The
universal property of profinite completion gives $c_{g,n}$, whose image
is closed and contains a dense subgroup, hence is the whole closure.
If its kernel is trivial, compactness makes it an isomorphism of
topological groups; pulling back open subgroups gives cofinality of the
congruence neighborhoods.  The converse follows from the same inverse
limit description.  A kernel calculation for a map into an ambient
group makes no assertion about conjugation by elements outside its
image.  Even for finite groups, an injective map $C_2\hookrightarrow
\mathfrak S_3$ need not have normal image.  Thus the replacement of all
the normality questions in scans 322--326 by the congruence problem, as
suggested in Lead 4, is invalid.
\end{proof}

\statement{B.2}{Known}{The established congruence cases}
For full profinite completion, $C_{g,n}=1$ when $g=0$, $n\geq3$; when
$g=1$, $n\geq1$; and when $g=2$, $n\geq0$.  The same faithfulness
holds for the finite extensions allowing permutations and orientation
reversal.  In particular, all genus-zero groups used in Part F have
their full profinite topology induced by their surface-group action.

The first two assertions are Asada's results, in the precise forms
recorded and proved in \citet[Proposition 2.6 and Theorem 2.8]{Boggi2009BC}.
The same attached paper proves the hyperelliptic result for punctured
surfaces, hence the genus-two result for $n\geq1$, in its Theorem 3.5.
The closed genus-two case requires the subsequent
\citet[Theorem 1.4(ii), with $\mathscr C=\mathscr F$]{Boggi2020BC}:
the hyperelliptic mapping class group equals the full mapping class
group in genus two.  The distinction matters: the 2009 paper expressly
restricts its main theorem to open surfaces and records a gap in an
earlier proof.  The more recent primary source
\citet[\S1.1 and Corollary B]{Boggi2026BC} confirms the genus-at-most-two
range and establishes additional centralizer cases, without solving
the problem in arbitrary genus.

To justify the finite-extension assertion, the pure orientation subgroup
is congruence-open: permutation of the finitely many peripheral classes
is continuous, and orientation reversal is distinguished from
orientation preservation by the multiplier modulo $3$.  Within the
geometric image the multiplier takes only the values $1$ and $-1$.
The kernel of the completed action on the finite extension therefore
lies in $\widehat P_{g,n}$, where it is trivial.  This is not the ordinary
congruence-subgroup property of $\operatorname{SL}_2(\mathbb Z)$ acting
on $\mathbb Z^2$.  For the punctured torus the decisive representation
acts on $\widehat F_2$ and remembers more than homology.

\statement{B.3}{Open}{The remaining full congruence problem}
For $g\geq3$, determine whether
\[
 \widehat P_{g,n}\longrightarrow
 \operatorname{Out}(\pi_{g,n})
\]
is injective.  The smallest genus and number of punctures not covered
by B.2 is $(g,n)=(3,0)$.  Part C proves
$C_{g,n+1}\simeq C_{g,n}$ whenever the target type is hyperbolic, so
adding punctures does not remove this unresolved case.  The exact
missing step is cofinality of characteristic-quotient levels among
all finite-index subgroups; neither the normality of a geometric image
in an arithmetic image nor faithfulness of the Galois action supplies
it.  This is a literature-open problem, not merely a gap in the
manuscript.  Our status check includes the 2026 primary source cited in
B.2.

\statement{B.4}{Proved}{Properness of the geometric image and the exceptional types}
For every hyperbolic type,
\[
 H_{g,n}\ \ne\ O_{g,n}.
\]
The assertion ``never an isomorphism'' in \S26, folio 163, scan 322
is correct in this range.  Read as an assertion about bare
fundamental groups outside that range, it is false; adding an
independent orientation line changes the data and must be specified.

\begin{proof}
The continuous multiplier
$\lambda:O_{g,n}\to\widehat{\mathbb Z}^{\times}$ takes values in
$\{1,-1\}$ on $H_{g,n}$.  Choose a smooth marked curve of the
specified type over some number field $L$; all required markings can
be made rational after a finite extension.  The arithmetic outer
action has multiplier the cyclotomic character of $G_L$.  Its image
is an open, hence infinite, subgroup of
$\widehat{\mathbb Z}^{\times}$, since
$[L\cap\mathbb Q(\mu_\infty):\mathbb Q]<\infty$.  Consequently an
arithmetic element with multiplier outside $\{1,-1\}$ lies in
$O_{g,n}\setminus H_{g,n}$.  No faithfulness theorem and no congruence
conjecture are needed for this argument.

For $(0,0)$ and $(0,1)$ the geometric fundamental group is trivial,
so its automorphism group and every image in it are trivial.  For
$(0,2)$ it is $\widehat{\mathbb Z}$, and the geometric image is
$\{1,-1\}\subset\widehat{\mathbb Z}^{\times}$; forgetting the two
labels does not change this image.  For $(1,0)$ it is
$\widehat{\mathbb Z}^2$, and the ambient group is
$\operatorname{GL}_2(\widehat{\mathbb Z})$.  The extended geometric
image is the subgroup with determinant in $\{1,-1\}$, whereas the
pure oriented image is $\operatorname{SL}_2(\widehat{\mathbb Z})$.
Here the action factors through the usual homological congruence
completion, and its kernel on the abstract profinite completion of
$\operatorname{SL}_2(\mathbb Z)$ is nontrivial.  Thus the punctured
torus and the unpunctured torus must not be merged.
The auxiliary orientation line that the manuscript adds in the
trivial-group cases, noted in A.11, is not contained in the bare
fundamental group; the preceding triviality statement does not
erase that extra datum.

For completeness, this last nontriviality has a finite witness.
The presentation $\operatorname{PSL}_2(\mathbb Z)=C_2*C_3$
maps onto $\operatorname{SL}_2(\mathbb F_8)$ by sending its
generators to
\[
 \begin{pmatrix}1&t\\0&1\end{pmatrix},\qquad
 \begin{pmatrix}0&1\\1&1\end{pmatrix},
 \quad\mathbb F_8=\mathbb F_2[t]/(t^3+t+1).
\]
The code verifies their orders, that they generate all 504
determinant-one matrices, and that the group is simple: its
conjugacy-class sizes are $1,56,56,56,56,63,72,72,72$, and no
proper union containing the identity has cardinality dividing
504.  A finite quotient of
$\operatorname{SL}_2(\widehat{\mathbb Z})$ has only cyclic
simple factors and groups $\operatorname{PSL}_2(\mathbb F_p)$
for primes $p\geq5$.  This follows by reduction modulo $p$
in each factor; the reduction kernel is pro-$p$, and the
groups for $p=2,3$ are solvable.  None of those projective
groups has order 504: the order formula gives 168 at $p=7$
and 660 at $p=11$, with no intervening prime.  The displayed
simple quotient therefore cannot be a congruence quotient,
which proves the claimed nontrivial kernel.
\end{proof}

\statement{B.5}{Refuted}{Normality of the full extended geometric group}
The assertion that $H_{g,n}$ is normal in $O_{g,n}$ for every
hyperbolic type is false already at $(0,3)$.

\begin{proof}
The three-punctured sphere has
$P_{0,3}=1$, $M^+_{0,3}=\mathfrak S_3$, and
$M^{\pm}_{0,3}=\mathfrak S_3\times C_2$.
For the rational curve
$\mathbb P^1\setminus\{0,1,\infty\}$ the outer Galois action is
faithful, as proved in D.16 from Belyi's theorem.  Its elements fix
the three peripheral labels.  The six rational fractional linear
transformations permuting $0,1,\infty$ commute with this action in the
outer automorphism group.  Hence
\[
 \mathfrak S_3\times G_{\mathbb Q}\hookrightarrow O_{0,3}.
\]
The intersection of the two factors is trivial because the first
factor acts faithfully on the labels and the second fixes them.
If $c\in G_{\mathbb Q}$ is complex conjugation, its image is the
orientation-reversing class fixing those labels.  Thus
$H_{0,3}=\mathfrak S_3\times\langle c\rangle$ inside this subgroup.
Normality of $H_{0,3}$ would imply normality of
$\langle c\rangle$ in $G_{\mathbb Q}$.

The splitting field of $t^3-2$ has Galois group $\mathfrak S_3$:
Eisenstein's criterion gives irreducibility and its discriminant is
$-108$, which is not a square.  Complex conjugation acts as a
transposition on its three roots.  Conjugating this transposition by
a three-cycle gives another transposition outside its order-two
subgroup.  Lifting the three-cycle to $G_{\mathbb Q}$ contradicts
normality.  The finite calculation, including the discriminant and
the conjugation, is reproduced by \texttt{code/bc\_counterexamples.py}.
\end{proof}

This answers the normality question attached to the extended group in
scan 322.  It does not answer the different question about
$\Sigma_{g,n}=\overline{M^+_{g,n}}$.  Nor does an author's annotation
``ultra-faux'' by itself provide a mathematical counterexample to a
different self-normalizer assertion: the conjecture so annotated is
recorded as withdrawn, and we do not silently substitute B.5 for it.

\statement{B.6}{Proved}{Conjugacy of geometric structures and their normalizers}
Fix an identification of a surface-group completion with a profinite
group with loops $\pi$, and let $\Sigma\subset O=\operatorname{Out}_{\mathrm{lac}}(\pi)$
be its oriented geometric image.  The conjugates of $\Sigma$ are
parameterized by $O/N_O(\Sigma)$.  The natural map
\[
 O/\Sigma\longrightarrow O/N_O(\Sigma)
\]
has fibers which are torsors under $N_O(\Sigma)/\Sigma$, after
transport by a representative.  In particular, uniqueness of this
conjugate subgroup means $N_O(\Sigma)=O$, equivalently normality of
$\Sigma$ in $O$.

\begin{proof}
Two conjugates $a\Sigma a^{-1}$ and $b\Sigma b^{-1}$ agree exactly
when $b^{-1}a$ normalizes $\Sigma$.  For a fixed representative $a$,
the fiber over $aN_O(\Sigma)$ is $aN_O(\Sigma)/\Sigma$.
This is the asserted torsor.  All subgroups used for the displayed
quotients are closed: the normalizer of a closed subgroup of a
profinite group is closed.  These are therefore profinite coset spaces,
although the first need not be a group.  This proves the assertions
in \S26, folios 164--165, scans 324--326, with the normalizer definition
kept explicit.
\end{proof}

\statement{B.7}{Open}{The unrestricted normalizer comparison}
The unrestricted questions
\[
 \Sigma_{g,n}\triangleleft O_{g,n},\qquad
 N^{\mathrm{norm}}_{g,n}=O_{g,n}
\]
are equivalent to each other by B.6 and are separate from B.3.
We do not resolve this unrestricted comparison.  At the smallest
test type $(0,3)$ it asks whether every inertia-preserving outer
automorphism normalizes the geometric $\mathfrak S_3$.
E.7 identifies the normalizer quotient using the additional
symmetry condition; that identification does not show that every
element of $O_{0,3}$ satisfies the condition.  Our searches of the
primary sources on quasi-special symmetric automorphisms and
procongruence automorphisms did not supply either this implication or
a counterexample.  Here the label Open records the precise unresolved
comparison in this audit; unlike B.3, we do not assert that each
individual case has been proved open by a literature classification.
The earlier self-normalizer conjecture for $H$, explicitly withdrawn
in scan 322, is not used as a premise.

\subsection{Actual discretizations under a finite covering}

A based discretization of $\pi$ will mean a dense subgroup
$D\subset\pi$, equipped with the oriented surface peripheral
structure, such that completion induces an isomorphism
$\widehat D\simeq\pi$ preserving that structure.  In particular it is
more than an arbitrary dense subgroup.  Forgetting orientation gives
the corresponding unoriented notion.  Quotienting by inner
automorphisms of $\pi$ gives strict prediscretizations; quotienting a
marking torsor by the closure of the geometric automorphism group
gives the manuscript's non-strict prediscretizations.  The adjective
``strict'' must not be dropped in these comparisons.

\statement{B.8}{Proved}{A surface lattice is self-normalizing in its completion}
If $D$ is a nonabelian free group or the fundamental group of a closed
orientable surface of genus at least two, then
\[
 N_{\widehat D}(D)=D.
\]

\begin{proof}
Free groups and closed orientable surface groups are conjugacy
separable and have Grossman's Property A: every automorphism taking
each element to a conjugate of itself is inner.  For the exact
conjugacy-separability statements we use
\citet[Theorems 3.2 and 4.1; Property A is stated in \S1 immediately
before Theorem 1.7]{Paris2009BC}.  These theorems are stronger, giving
conjugacy $p$-separability.  If $h\in\widehat D$ normalizes $D$,
conjugation by $h$ induces an automorphism $\alpha$ of $D$.  For every
$d\in D$, $\alpha(d)$ and $d$ are conjugate in $\widehat D$; hence
they are conjugate in $D$.  Property A gives
$\alpha=\operatorname{int}(d_0)$ for some $d_0\in D$.
Then $d_0^{-1}h$ centralizes the dense subgroup $D$ and therefore
centralizes $\widehat D$.  The latter is center-free.  For example,
this follows from the slimness theorem in D.14, or
from the centralizer calculation for two independent simple loops
in \citet[Proposition 3.6]{BoggiZalesskii2019BC}.
Consequently $h=d_0$.  The reverse inclusion is immediate.
\end{proof}

\statement{B.9}{Proved}{Injectivity of restriction to a finite cover}
Let $\pi$ have hyperbolic orientable surface type and let
$V\subset\pi$ be an open subgroup with its induced peripheral
structure.  The map
\[
 \operatorname{Discr}(\pi)\longrightarrow\operatorname{Discr}(V),
 \qquad D\longmapsto D\cap V
\]
is well defined and injective, for based discretizations, and the
same holds with orientation retained.  No congruence-subgroup
property is assumed.  This proves the corollary in \S28, folio 178,
scans 352--353.

\begin{proof}
The intersection $D\cap V$ is of finite index in $D$.  The topology
induced on a finite-index subgroup by the full profinite topology is
its full profinite topology: the core in $D$ of a subgroup of finite
index in $D\cap V$ is of finite index in $D$.  Therefore
$\widehat{D\cap V}\simeq V$.  The ordinary finite covering of
surfaces determines the asserted peripheral and orientation data.

First suppose that $V$ is normal.  If
$D_1\cap V=D_2\cap V=\Lambda$, both $D_1$ and $D_2$ normalize
$\Lambda$.  B.8 applied inside $V=\widehat\Lambda$ gives
\[
 N_\pi(\Lambda)\cap V=\Lambda.
\]
Each $D_i$ maps onto the finite group $\pi/V$, since it is dense.
For $a\in N_\pi(\Lambda)$ choose $d\in D_i$ with the same image
in $\pi/V$.  Then $ad^{-1}\in\Lambda\subset D_i$, proving
$N_\pi(\Lambda)=D_i$.  Thus $D_1=D_2$.
For arbitrary $V$, replace it by its core
$W=\bigcap_{a\in\pi}aVa^{-1}$.  This is open and normal, and
$D_1\cap W=D_2\cap W$ follows by intersecting the given equality
with $W$.  The normal case applies.  Orientation is unchanged under
these restrictions.
\end{proof}

\statement{B.10}{Proved}{The image criterion for a normal finite cover}
Suppose $V\triangleleft\pi$ is open, $G=\pi/V$, and
$\Lambda\subset V$ is an oriented discretization.  It extends to an
oriented discretization of $\pi$ if and only if the faithful deck
outer action
\[
 G\longrightarrow\operatorname{Out}_{\mathrm{lac}}(V)
\]
factors through the image of
$\operatorname{Out}_{\mathrm{lac}}^+(\Lambda)$ in
$\operatorname{Out}_{\mathrm{lac}}(V)$.  The extension, when it
exists, is unique as a subgroup of $\pi$.

\begin{proof}
Necessity is immediate by conjugating $D\cap V$ inside an extending
lattice $D$.  For sufficiency, the discrete outer action on $\Lambda$
is unique because the completion map on its outer automorphism group
is injective, by the same conjugacy-separability and Property A
argument as B.8.  The deck action is faithful: an element of $\pi$
acting on $V$ by an inner automorphism differs from an element of
$V$ by an element centralizing $V$, and $C_\pi(V)=1$.  The last
equality is the slimness of a hyperbolic profinite surface group.

Since $Z(\Lambda)=1$, the specified outer action determines a unique
extension $1\to\Lambda\to E\to G\to1$.  Explicitly, it is the
pullback of $\operatorname{Aut}(\Lambda)\to
\operatorname{Out}(\Lambda)$ along that action.  Completing
automorphisms maps it into the analogous pullback for $V$; the
latter is $\pi$, since $Z(V)=1$.  The resulting map $E\to\pi$ is
injective, induces the given inclusion on $\Lambda$, and is the
identity on $G$.  Its image is dense.  Because $\Lambda$ has finite
index in $E$, completion of this extension is exact, and the
comparison with $1\to V\to\pi\to G\to1$ gives
$\widehat E\simeq\pi$.

The profinite group $\pi$ is torsion-free.  A precise source is
\citet[Lemma 3.1(ii)]{BoggiZalesskii2019BC}, applied to a Fuchsian
surface group without orbifold points: there are no decomposition
groups in which a nontrivial finite subgroup could lie.
Thus $E$ is torsion-free.  The finite extension is allowable in the
Fuchsian sense: its outer action preserves the peripheral classes
and orientation.  \citet[Theorem 7, p.~258]{Kerckhoff1983BC} says
that an allowable finite extension of a Fuchsian group is Fuchsian
when its outer action is injective.  Hence $E$ is the fundamental
group of an orientable surface of finite type, and the covering
defined by $\Lambda$ is a genuine unramified surface covering.

It remains to check the loops, rather than merely the underlying
group.  Each peripheral cyclic subgroup of $E$ meets $\Lambda$ in
a finite-index peripheral subgroup.  Its closure is the centralizer
in $\pi$ of the corresponding peripheral subgroup of $V$.
Conversely, the original peripheral groups of $\pi$ are obtained
from those of $V$ by exactly this centralizer operation.
Both assertions follow from
\citet[Proposition 3.6]{BoggiZalesskii2019BC}: a nonzero power of
an algebraically simple element has normalizer, and hence
centralizer, equal to the maximal procyclic group it generates.
The two peripheral structures therefore agree.  The allowable
orientation agrees on the finite cover and hence on the base.
Uniqueness is B.9.
\end{proof}

In scan 352 the displayed quotient for the deck group is transcribed
as $\pi'/\pi$, although the surrounding text has $\pi'\subset\pi$.
The image inspection confirms this apparent slip.  B.10 requires
$G=\pi/\pi'$; it does not assert that the handwriting has thereby
been editorially settled.  The same page only sketches the
loop-compatibility step, supplied above using the centralizer theorem.

\statement{B.11}{Refuted}{Surjectivity for raw discretizations of a finite cover}
Restriction of based discretizations need not be surjective.  A
counterexample is the degree-two covering
\[
 \mathbb P^1(\mathbb C)\setminus\{0,\infty,1,-1\}
 \longrightarrow
 \mathbb P^1(\mathbb C)\setminus\{0,\infty,1\},
 \qquad z\longmapsto z^2.
\]

\begin{proof}
Let $V=\widehat\Lambda$ be the subgroup defined by this cover in
$\pi=\widehat F_2$, and let $s$ be the deck involution in
$\operatorname{Out}_{\mathrm{lac}}(V)$.  It fixes the punctures
$0,\infty$ and interchanges $1,-1$.  In the discrete pure mapping
class group $P_{0,4}\simeq F_2$, conjugation by $s$ is nontrivial:
it takes the twist about a curve separating $\{0,1\}$ from the
other two punctures to the distinct twist about a curve separating
$\{0,-1\}$.  The permutation of these two partitions is checked in
\texttt{code/bc\_counterexamples.py}.

The centralizer of $s$ in $P_{0,4}$ has infinite index.  Otherwise
the induced automorphism would fix a subgroup of finite index.
Each element of $F_2$ has a positive power in that subgroup, and
the unique-root property in a free group would force the
automorphism to fix every element, a contradiction.  Let
$B=\overline{P_{0,4}}\subset\operatorname{Out}_{\mathrm{lac}}(V)$.
Its closed subgroup $C_B(s)$ is not open, since openness would
make its intersection with the dense group $P_{0,4}$ of finite
index.  Thus $B/C_B(s)$ is uncountable: a compact Hausdorff space
cannot be a countable union of closed cosets with empty interior.

The discrete group $\operatorname{Out}_{\mathrm{lac}}^+(\Lambda)$
is countable.  We may therefore choose $u\in B$ for which
$u^{-1}su$ is not in that discrete group.  For definiteness this
choice can be specified by enumerating the discrete group and
choosing successively nested basic clopen neighborhoods which
avoid the closed conjugator set for the next element; require
also containment in the next characteristic finite-quotient
neighborhood.  Each avoidance is possible because the nonempty
conjugator sets are cosets of the nonopen centralizer.  Compactness
and the cofinal neighborhood condition determine an element $u$.
Choose a genuine automorphism representing $u$, and put
$\Lambda'=u(\Lambda)$.  The deck outer action does not factor
through $\operatorname{Out}_{\mathrm{lac}}^+(\Lambda')$.
By B.10, $\Lambda'$ does not extend to $\pi$.

The counterexample is necessarily profinite: all discretizations
are dense and their difference cannot be certified by distinguishing
their images in a single finite quotient.  The nested construction
above specifies the required element, rather than relying on a
purported finite-quotient witness.  Notice also that $u\in B$:
$\Lambda$ and $\Lambda'$ have the same non-strict oriented
prediscretization.  This counterexample does not disprove the
prediscretization assertion in scan 355.
\end{proof}

\subsection{Passing to prediscretizations and the remaining comparisons}

\statement{B.12}{Proved}{Restriction descends to both prediscretization quotients}
For every finite covering $V\subset\pi$ as in B.9, restriction of
discretizations descends to strict prediscretizations and to the
non-strict oriented and unoriented prediscretizations.  In
particular, the well-definedness doubts in \S28, folios 179--181,
scans 355--358 have an affirmative answer.

\begin{proof}
First let $D_2=aD_1a^{-1}$, $a\in\pi$.  Density gives
$\pi=VD_1$: every coset of the open subgroup $V$ meets $D_1$.
Write $a=vd$ with $v\in V$, $d\in D_1$.  Then
$D_2=vD_1v^{-1}$, and
$D_2\cap V=v(D_1\cap V)v^{-1}$.  This is precisely the correction
already made in scan 357.

For the closure quotient, fix $D=D_1$ and let $A_D$ be its discrete
loop-preserving automorphism group, with orientation restriction
if required, inside $\operatorname{Aut}_{\mathrm{lac}}(\pi)$.
Suppose $D_2=u(D)$ with $u\in\overline{A_D}$.  The group $\pi$ is
topologically finitely generated, so it has finitely many open
subgroups of index $[\pi:V]$.  The action on this finite set is
continuous.  Density of $A_D$ in its closure permits a choice
$a\in A_D$ with $a(V)=u^{-1}(V)$.  Thus $ua$ stabilizes $V$ and
$ua(D)=D_2$.

The intersection of $\overline{A_D}$ with the open stabilizer of
$V$ is the closure of $A_D\cap\operatorname{Stab}(V)$.
Restriction of automorphisms stabilizing $V$ to $V$ is continuous;
for instance, restriction is a continuous map between the closed
automorphism spaces in the spaces of maps on the corresponding
finitely generated profinite groups.  It takes the discrete
stabilizer into $\operatorname{Aut}_{\mathrm{lac}}(D\cap V)$.
Therefore the restrictions of $D$ and $u(D)$ differ by an element
in the closure of that latter group.  The same proof preserves
orientation, because orientation restricts compatibly along a
surface covering.  It proves descent, not bijectivity.
\end{proof}

\statement{B.13}{Proved}{The exact coset criteria behind conditions (a)--(e)}
Let $f:A\to A'$ be a continuous homomorphism of profinite groups,
and let $B\subset A$, $B'\subset A'$ be closed subgroups with
$f(B)\subset B'$.  The map of right-coset spaces
$A/B\to A'/B'$ is injective if and only if
\[
 f^{-1}(B')=B,
\]
and is surjective if and only if $A'=f(A)B'$.  If in addition
$f$ is surjective and $\ker(f)\subset B$, then
\[
 B=f^{-1}(B')\quad\Longrightarrow\quad
 N_A(B)=f^{-1}(N_{A'}(B')),
 \qquad N_A(B)/B\simeq N_{A'}(B')/B'.
\]

\begin{proof}
The fiber over the identity coset is $f^{-1}(B')/B$, and translating
by an element of $A$ gives all other nonempty fibers.  The image
consists exactly of cosets meeting $f(A)$.  For the last assertion,
if $B=f^{-1}(B')$ and $f$ is onto, an element conjugates $B$ onto
itself exactly when its image conjugates $B'$ onto itself.
The map on normalizers is therefore onto and its kernel is
$\ker(f)\subset B$; quotienting gives the isomorphism.
This proves the precise formal mechanism, and also specifies the
extra hypotheses, behind \S28, folio 181, scans 358--359.  In a
restriction map to a characteristic open subgroup one need not
have surjectivity of $f$; the displayed last assertion cannot be
used without verifying it or an appropriate replacement.
\end{proof}

\statement{B.14}{Open}{Bijectivity for arbitrary profinite prediscretizations}
For a characteristic open subgroup $V\subset\pi$, put
\[
 A=\operatorname{Aut}_{\mathrm{lac}}(\pi),\quad
 A'=\operatorname{Aut}_{\mathrm{lac}}(V),\quad
 B=\overline{\operatorname{Aut}_{\mathrm{lac}}^+(D)},\quad
 B'=\overline{\operatorname{Aut}_{\mathrm{lac}}^+(D\cap V)}.
\]
Restriction gives a continuous injection $f:A\to A'$, and B.12
gives $f(B)\subset B'$.  The unrestricted bijectivity proposed in
scan 355 requires both
\[
 f^{-1}(B')=B,\qquad A'=f(A)B'.
\]
We have not proved these two assertions for arbitrary
loop-preserving automorphisms.  A concrete characteristic test,
already indicated in \S28, folio 183, scan 363, is
\[
 \pi=\widehat F_2,\qquad
 V=\ker(\widehat F_2\to(\mathbb Z/2)^2),
\]
the degree-four covering of the three-punctured sphere of type
$(0,6)$.  Among noncharacteristic coverings, the degree-two cover
in B.11 is the first basic test of the set map of B.12; its raw
failure of surjectivity disappears in the closure quotient and
does not decide this test.  The congruence property is known for
both surfaces in these tests, so restating these conditions as
CSP does not supply the missing argument.  The analogous filling
conditions are made exact in C.8.  As in B.7, Open here records
the unresolved unrestricted statement in this audit, without
claiming a literature classification of every individual cover.
For completeness, injectivity of the restriction homomorphism
uses slimness, not CSP: if an automorphism is the identity on
$V$, then for each $a\in\pi$ the element $a^{-1}f_0(a)$
centralizes the normal open subgroup $V$, and hence is trivial
by D.14--D.15.  Here $f_0$ denotes that automorphism, to
distinguish it from the restriction homomorphism $f$.

\statement{B.15}{Refuted}{Raw discretization injectivity does not transfer to filling}
If the proposed injection of raw discretizations in item 3 of the
prompt is applied also to filling a puncture, it is false.  It
already fails for $(0,4)\to(0,3)$.

\begin{proof}
Let $D=F(x,y,z)$ be the fundamental group of a four-punctured
sphere with peripheral representatives $x,y,z,(xyz)^{-1}$,
and let $\pi=\widehat D$.  Filling the puncture represented by
$x$ gives
\[
 q:\pi\longrightarrow\widehat F(y,z).
\]
Choose $\alpha\in\widehat{\mathbb Z}=\prod_p\mathbb Z_p$ with
$\alpha_2=0$ and $\alpha_p=1$ for every odd prime.  This is not
an integer: an integer equal to zero in $\mathbb Z_2$ is zero,
whereas its $3$-adic component here is one.  Set $h=x^\alpha$.
The retraction $D\to\langle x\rangle$, extended to completions,
shows that
$D\cap\overline{\langle x\rangle}=\langle x\rangle$;
hence $h\notin D$.  B.8 gives $hDh^{-1}\ne D$.
Both are oriented discretizations, because inner automorphisms
preserve all peripheral conjugacy and orientation data.  But
$q(h)=1$, so
\[
 q(hDh^{-1})=q(D)=F(y,z).
\]
Thus the raw filling map is not injective.  The code records
compatible finite Chinese-remainder approximations to $\alpha$;
the proof of nonintegrality is the argument just given, not a
claim that finitely many residues certify it.  The two
discretizations are inner-conjugate, so this counterexample leaves
both versions of the prediscretization question untouched.
\end{proof}

All statements in this part concern orientable surfaces with the
specified orientation option.  Allowing orientation reversal does
not turn them into statements about fundamental groups of
nonorientable surfaces.  A real curve is handled by its arithmetic
extension over $\mathbb R$, whose quotient is $C_2$; the image of
complex conjugation is exactly why the full extended group in
B.5 behaves differently from the oriented geometric image.  The
full profinite results retain the primes $2$ and $3$; no argument
has inverted either prime.  The use of reduction modulo $3$ in
B.2 only distinguishes the two orientation signs.
